% blog.tex

\newcommand{\BlogTitle}{Blog techniczny}
\newcommand{\BlogNote}{}
\renewcommand{\BlogList}{}
\AddBlogEntry
  {Analiza wydajności arytmetyki JVM z perspektywy kodu bajtowego oraz instrukcji maszynowych x86}
  {25 sty 2026}
  {%
    Studium przypadku weryfikujące mechanizmy "strength reduction" w maszynie wirtualnej HotSpot. Analiza obejmująca
    transformację arytmetyki kodu bajtowego na natywne przesunięcia bitowe oraz operacje logiczne procesora x86.
  }
  {analiza-wydajności-arytmetyki-jvm-z-perspektywy-kodu-bajtowego-oraz-instrukcji-maszynowych-x86}
\AddBlogEntry
  {Zespolona interpretacja równania różniczkowego w ujęciu transformacji Fouriera}
  {15 gru 2025}
  {%
    O tym, jak wzór Eulera zamienia równania różniczkowe na prostą impedancję cewki oraz dlaczego splot w dziedzinie
    czasu to proste mnożenie w analizie Fouriera.
  }
  {zespolona-interpretacja-równania-różniczkowego-w-ujęciu-transformacji-fouriera}
\AddBlogEntry
  {Trochę o JWT i dlaczego nie jest taki świetny, jak się na pozór wydaje}
  {20 lis 2025}
  {O minusach technologii JWT i dlaczego czasami warto postawić na starą dobrą stanową sesję.}
  {trochę-o-jwt-i-dlaczego-nie-jest-taki-świetny-jak-się-na-pozór-wydaje}
\AddBlogEntry
  {Szyfrowanie obrazów z wykorzystaniem chaotycznego atraktora Lorenza}
  {22 cze 2025}
  {Koncepcyjny sposób szyfrowania strumieniowego z wykorzystaniem punktów przestrzeni fazowej atraktora.}
  {szyfrowanie-obrazów-z-wykorzystaniem-chaotycznego-atraktora-lorenza}
\AddBlogEntry
  {Prosty, bazujący na UDP i AES szyfrowany protokół transmisji obrazu}
  {28 gru 2024}
  {Opis zaprojektowanego i zaimplementowanego przeze mnie szyfrowanego protokołu transmisji obrazu.}
  {prosty-bazujący-na-udp-i-aes-szyfrowany-protokół-transmisji-obrazu}
\AddBlogEntry
  {Podstawy przetwarzania i analizy częstotliwościowej sygnału audio}
  {20 kwi 2024}
  {%
    Techniczne podstawy przetwarzania sygnału audio, reprezentacja matematyczna, funkcje sinusoidalne i piłokształtne
    oraz analiza utworu muzycznego.
  }
  {podstawy-przetwarzania-i-analizy-częstotliwościowej-sygnału-audio}

% education.tex

\newcommand{\EducationTitle}{Edukacja}
\newcommand{\ThesisLabel}{Praca dyplomowa}
\renewcommand{\EducationList}{}
\AddEducationEntry
  {Politechnika Śląska \& NASK}
  {Gliwice, Polska}
  {Studia podyplomowe -- Cyber Science (Zarządzanie Cyberbezpieczeństwem)}
  {2025 -- obecnie}
  {%
    Zaprojektowanie kryptograficznej architektury Zero Trust zgodnej z dyrektywą NIS2, wykorzystującej HashiCorp Vault,
    post-kwantowe mTLS (ML-KEM/ML-DSA) oraz autorskie szyfrowanie na poziomie aplikacji w wysoko dostępnym środowisku
    Kubernetes.
  }
\AddEducationEntry
  {Politechnika Śląska}
  {Gliwice, Polska}
  {Magister inżynier informatyki; ocena: 5.0 (4.82)}
  {2024 -- 2025}
  {%
    Analiza pozbawionej podpróbkowania diadycznego dyskretnej transformacji falkowej, wzbogaconej o nowatorski estymator
    wag podpasm, w celu detekcji krawędzi na obrazach zniekształconych wysokoczęstotliwościowym szumem impulsowym.
  }
\AddEducationEntry
  {Politechnika Śląska}
  {Gliwice, Polska}
  {Inżynier informatyki; ocena: 5.0 (4.67)}
  {2020 -- 2024}
  {%
    Demonstracja możliwości komunikacyjnych w silnie rozproszonej infrastrukturze mikroserwisowej z wykorzystaniem
    architektury sterowanej zdarzeniami Apache Kafka w komunikacji czasu rzeczywistego (aplikacja typu chat).
  }

% experience.tex

\newcommand{\ExperienceTitle}{Doświadczenie zawodowe}
\newcommand{\WorkHybrid}{hybrydowo}
\renewcommand{\ExperienceList}{}
\AddExperienceEntry
  {Euvic}
  {Gliwice, Polska (\WorkHybrid)}
  {Stażysta Full Stack Developer}
  {lip 2023 -- paź 2023}
  {Next.js, Nest.js, PostgreSQL, MongoDB, Docker, AWS Cognito, Agile, Jira, współpraca z zespołem z UK}

% header.tex

\newcommand{\LabelEmail}{E-mail:}
\newcommand{\LabelPortfolio}{Portfolio:}
\newcommand{\LabelMobile}{Telefon:}
\newcommand{\LabelGitHub}{GitHub:}
\newcommand{\LabelLinkedIn}{LinkedIn:}

% hobby.tex

\newcommand{\TitleHobby}{Zainteresowania}
\newcommand{\TextHobby}{
  Entuzjasta systemów automatyki domowej i multimedialnej. Archeolog internetowy, miłośnik kina studyjnego i
  alternatywnych brzmień. Aktywny archiwista i użytkownik nośników analogowych. Czytelnik literatury science fiction,
  kosmologicznej i Lovecraft'owskiego horroru.
}

% honors_and_awards.tex

\newcommand{\AwardsTitle}{Nagrody i wyróżnienia}
\renewcommand{\AwardsList}{}
\AddAwardEntry
  {Stypendium Rektora za wybitne osiągnięcia naukowe}
  {paź 2024}
  {Politechnika Śląska -- znalezienie się w gronie 6\% najlepszych studentów na kierunku.}

% projects.tex

\newcommand{\ProjectsTitle}{Wybrane projekty}

% research.tex

\newcommand{\IndependentResearchTitle}{Niezależna działaność badawcza}
\AddResearchEntry
  {Steganografia w dziedzinie transformat falkowych}
  {2026 -- obecnie}
  {
    Badania nad wykorzystaniem stałoprzecinkowej DTCWT, modulacji QAM oraz technik multimodalnych (detekcja transientów)
    w celu ukrywania danych w sygnałach obrazowych.
  }
\AddResearchEntry
  {Estymacja wag podpasm dekompozycji falkowej}
  {2025}
  {
    Model redukcji szumu, w którym aproksymowana funkcja Gaussa wyznaczała wagi kolejnych podpasm dekompozycji falkowej.
    Rozwiązanie to posłużyło jako podstawa pracy magisterskiej.
  }

% rodo.tex

\newcommand{\TextRodo}{
  Wyrażam zgodę na przetwarzanie moich danych osobowych dla potrzeb niezbędnych do realizacji procesu rekrutacji
  (zgodnie z ustawą z dnia 10 maja 2018 roku o ochronie danych osobowych (Dz. U. z 2018, poz. 1000) oraz zgodnie z
  Rozporządzeniem Parlamentu Europejskiego i Rady (UE) 2016/679 z dnia 27 kwietnia 2016 r. w sprawie ochrony osób
  fizycznych w związku z przetwarzaniem danych osobowych i w sprawie swobodnego przepływu takich danych oraz uchylenia
  dyrektywy 95/46/WE (RODO).
}

% skills.tex

\newcommand{\SkillsTitle}{Technologie i umiejętności}
\renewcommand{\SkillsList}{}
\AddSkillEntry
  {Architektura}
  {
    mikroserwisy, rozpraszanie usług stanowych, architektura hexagonalna (porty i adaptery), projektowanie API i
    bibliotek (DI, IoC, silniki oparte o adnotacje), wzorce projektowe, SOLID
  }
\AddSkillEntry
  {Testowanie}
  {JUnit, Mockito, AssertJ, Testcontainers (testy integracyjne), TDD}
\AddSkillEntry
  {Backend}
  {
    JVM (Java, Kotlin), Spring Framework, Spring Boot, Maven, Gradle, Jetty, GraalVM (polyglot), RabbitMQ, JNI, Apache
    Kafka
  }
\AddSkillEntry
  {Protokoły}
  {TCP, UDP, HTTP, REST, WebSocket, Protobuf, MQTT (Moquette), mDNS}
\AddSkillEntry
  {Bazy danych}
  {MySQL, PostgreSQL, SQLite, Redis (KV, Pub/Sub), Cassandra, Neo4j}
\AddSkillEntry
  {Systemy wbudowane (embedded)}
  {C, CMake, ESP-IDF (FreeRTOS), ESP32/ESP32-C3 (RISC-V), I2C, SPI, RS-232/485, aktualizacje OTA}
\AddSkillEntry
  {R\&D i DSP}
  {Python, OpenCV, PyWavelets, NumPy, SciPy, Pandas, \LaTeX}
\AddSkillEntry
  {Frontend, mobile i desktop}
  {
    TypeScript, JavaScript, Next.js, React (Redux Toolkit, TanStack Query), Android (Kotlin, Jetpack Compose), Java
    Swing, Electron
  }
\AddSkillEntry
  {Chmura i DevOps}
  {
    Git (GitOps), Docker, GitHub Actions, Linux (Debian), Z shell, GCP, OCI, Nginx, HashiCorp Consul (magazyn KV),
    Prometheus, Grafana, Loki, K8S
  }
\AddSkillEntry
  {Bezpieczeństwo}
  {HashiCorp Vault, Cloudflare Zero Trust}
